\chapter{The host organisation}
\label{ch:host}

\section{Forschungszentrum Jülich}

Forschungszentrum Jülich GmbH is located between Aachen and Cologne, in North
Rhine-Westphalia, and is among the largest interdisciplinary research
institutions in Europe. It is a member of the Helmholtz Association, Germany's
largest scientific organisation, whose research programmes are organised
around problems that require infrastructure and continuity beyond what a
single university group typically sustains.
\needcheck{Confirm current staff and budget figures against the official FZJ
institutional profile before submission.}

The centre runs research programmes spanning energy, information technology,
and the bioeconomy, alongside the physical and computational infrastructure
these programmes require. This organisational structure was directly relevant
to the present work: the project combines computer vision and geometry
processing methods with morphometric analysis intended for use by biologists,
and was executed in part on a national supercomputing facility located on the
same campus as the biological collections the resulting measurements are
intended to serve. Few research environments make this combination of
disciplines and infrastructure available within a single site.

From the perspective of the internship, the most important feature of the host
institution is not only its scientific breadth but also its ability to connect
research questions to the infrastructure required to answer them. The project
involved scanning, reconstruction, geometry processing, optimisation,
validation, and downstream morphometric interpretation. Each of these stages
depends on a different technical ecosystem, and the institute provides a setting
where those ecosystems coexist under one organisation rather than being spread
across disconnected partners. This environment was crucial for the method
development reported here, because the work was not limited to algorithm design
in isolation; it required access to large-scale computation, scientific data
handling, and active collaboration across disciplinary boundaries.

\subsection{The Jülich Supercomputing Centre}

The Jülich Supercomputing Centre (JSC) is one of Europe's principal
high-performance computing facilities, hosting systems in the \term{jureca} and
\term{juwels} families.
\needcheck{Confirm current JSC system lineup and any exascale system status at
time of submission.}

The experimental work reported here was executed on \term{jureca}: the software
environment was built under Conda on the scratch filesystem, jobs were
dispatched to GPU-equipped compute nodes, and long-running fits were detached
under \code{tmux} to allow execution to continue independently of an active
terminal session. Several of the experiments reported in later chapters, in
particular the scale-barrier validation at $N=838$ specimens and the
multi-seed benchmarking described in Chapter~\ref{ch:elimination}, require
compute resources beyond what a single workstation can provide within the
timeframe of the internship; their inclusion in this report depended on access
to this infrastructure.

Working with the cluster also required adapting standard practice for
reproducible, long-running computation: environment management under Conda,
job scheduling, checkpointing of intermediate results, and a clear separation
between login nodes, which are not GPU-equipped and are intended for
compilation and light interactive work, and compute nodes, on which numerical
work is actually executed. This distinction was not initially obvious and
its absence was the direct cause of several early failed job submissions; it
is recorded here because it materially affected how many experiments could be
run within the available time, and therefore how much could be established
empirically over the course of the internship.

\section{Research context: parametric models for non-model species}

The immediate scientific context for this internship is a line of work
concerned with extending parametric three-dimensional shape models to species
outside the small set that computer vision research has historically
addressed.

Behavioural and biomechanical analysis increasingly depends on recovering
the pose and shape of an animal from imperfect observations. The corresponding
tooling is well developed for humans and largely unavailable elsewhere.
\term{smpl}~\cite{smpl} and its successors assume a cooperative subject, a
large scan corpus, and a fixed, well-understood body topology. \term{smal}~%
\cite{smal} relaxed the first two of these assumptions for quadrupeds by
learning a shape space from toy figurines rather than live animals, and
demonstrated that a model trained in this way generalises to photographs of
real animals, including species not present during training.

\term{smilify} extends this line of work to insects, where the body plan
itself differs substantially from the mammalian case and the anatomical
structures of interest are thin, numerous, and highly articulated. Upstream of
the registration pipeline is \term{scAnt}~\cite{scant}, an open-source macro
photogrammetry scanner designed to make high-resolution three-dimensional
capture of small specimens accessible to laboratories without dedicated
scanning infrastructure. The two systems are designed to complement one
another: an affordable scanner produces a large volume of imperfect scan data,
and a robust registration pipeline is what converts that data into comparable,
quantitative measurements. Neither component is of substantial independent
value without the other.

The broader objective underlying both systems, and the reason this project was
of particular interest, is the extension of quantitative shape and pose
analysis to research contexts that cannot support the data collection
requirements of the human-centred literature: specifically, laboratories
studying non-model species, using affordable equipment, on data that reflects
the practical limitations of real-world scanning rather than idealised
laboratory conditions.

\section{Supervision and working arrangement}

The internship was hosted by Dr.\ Fabian Plum, postdoctoral researcher and
author of both \term{scAnt} and \term{smilify}, and was conducted on site from
1 July to 25 September 2026.
\needcheck{Confirm the exact institute or division name within FZJ for the
title page.}

Two aspects of the working arrangement are relevant to how this report should
be read. First, the assignments listed in the technical data sheet were
presented as an initial hypothesis about where effort would be most valuable,
rather than as a fixed specification. When the controlled comparison described
in \S\ref{sec:turn} indicated that mesh input quality was not the limiting
factor, redirecting the remaining work toward initialisation and
correspondence was treated as the correct response to the evidence, consistent
with the framing under which the internship was supervised.

Second, the work was conducted on a live, shared repository, with code review
conducted through standard pull-request practice and with the supervisor
himself running related experiments on adjacent branches. \term{smilify}
supports authoring per-joint rotation limits directly in Blender, using
standard \emph{Limit Rotation} bone constraints on the template armature;
these are exported into the model file as a \code{joint\_limits} array and
were, at the start of this internship, consumed as a pose prior by the
neural-inference fitting path but not by the 3D mesh-registration path in
\code{fitter\_3d/}. This gap was identified in supervisor review, scoped as
GitHub Issue~\#97, and closed as part of this internship by adding the
corresponding hinge-loss term to the registration optimiser, so that limits
authored once in Blender are now honoured consistently across all of
\term{smilify}'s fitting paths.

Several results reported in Chapter~\ref{ch:elimination} were revised
following supervisor review, including corrections to aggregate statistics
and to an initially incorrect attribution of a failure mode; these corrections
are noted at the relevant points in that chapter, since the pre-registration
discipline used throughout this report was adopted in direct response to that
review process.

Beyond ants, the same registration and neural-inference pipeline has been used
within the project to build parametric models of other species, including a
mouse model and a multi-species stick-insect model trained on data from an
omnidirectional treadmill. This breadth is relevant context for the present
work: the correspondence and initialisation problems investigated here are not
specific to Formicidae, and a fix validated on ants is, in principle, a fix
available to every species the framework is extended to next.

\section{Academic framing}

This internship forms part of the second year of the engineering cycle at
Université Mohammed VI Polytechnique, and the present report is the
corresponding academic deliverable. Its purpose is to document both the
technical contributions produced during the internship and the evidentiary
basis for the decisions that shaped the direction of the work, since in this
project the two are not separable: the most consequential decisions taken
during the internship concerned which lines of investigation to discontinue,
and these decisions are only defensible if the evidence behind them is
presented alongside the results.
